\section{Function Study}

\begin{exercise}
\index{problem!of the can}%
\index{can!problem of the}%
Determine the radius and height of a cylindrical can with a volume of $33 cl$
that, for the same volume, has the minimum total surface area.
\end{exercise}
\begin{proof}[Solution.]
Let $V$ be the volume, $S$ the total surface area, $h$ the height, and $r$ the
base radius of the cylinder. We know that
\[
  V = \pi h r^2, \qquad
  S = 2\pi r h + 2 \pi r^2.
\]
Solving for $h$ from the first equation and substituting into the second, we
obtain
\[
  S = 2 \pi r \frac{V}{\pi r^2} + 2 \pi r^2
    = \frac{2V}{r} + 2 \pi r^2.
\]
The function $S(r)$ is defined and continuous on $(0,+\infty)$ and we have
$S(r)\to +\infty$ as $r\to 0^+$ and also as $r\to +\infty$. Thus, $S$ admits a
minimum by the generalized Weierstrass theorem. To find the minimum, it suffices
to compute the derivative
\[
 \frac{dS}{dr} = -\frac{2V}{r^2} + 4 \pi r = \frac{4\pi r^3 - 2V}{r^2}
\]
and find the critical points
\[
  4\pi r^3 = 2V
\]
from which
\begin{align*}
 r = \sqrt[3]{\frac{V}{2\pi}} \approx 3.74 cm\\
 h = \frac{V}{\pi r^2} \approx 7.49 cm.
\end{align*}
\end{proof}

\begin{exercise}
Solve the equation
\begin{equation} \label{eq:4734521}
  e^x = x^3.
\end{equation}
\end{exercise}
%
\begin{proof}[Solution.]
The left side of the equation is a positive number, and thus we can certainly
assume $x>0$; otherwise, the right side would not be positive as well. We can
then take the logarithm of both sides and obtain the equation
\[
  x = 3 \ln x
\]
which is equivalent to the given equation.

Consider the function
\[
 f(x) = x - 3 \ln x
\]
so that the solutions sought are the zeros of $f$. We have
\[
  f'(x) = 1 - \frac{3}{x}
\]
and we can therefore assert that
\mynote{The symbol $\gtrless$ is used to indicate that this inequality and the following ones can be written with either the $>$ or $<$ sign as long as the same sign is used throughout.}%
$f'(x) \gtrless 0$
if and only if
\[
  1 \gtrless \frac 3 x
\]
that is (recalling that we are assuming $x>0$)
\[
  x \gtrless 3.
\]
That is: if $x>3$, then $f'(x)>0$ and consequently
(theorem~\ref{th:criteri_monotonia}) $f$ is strictly increasing when restricted
to the interval $[3,+\infty)$, whereas if $x<3$, then $f'(x)<0$ and consequently
$f$ is strictly decreasing when restricted to the interval $(0,3]$. To determine
if the function $f$ vanishes at some point, we now need to evaluate $f$ at the
endpoints of the intervals on which it is monotonic. We have
\[
  \lim_{x\to 0^+} f(x) = +\infty, \qquad
  f(3) = 3 (1 - \ln 3) <0, \qquad
  \lim_{x\to +\infty} f(x) = +\infty.
\]
We observe that at the endpoints of the intervals $(0,3]$ and $[3,+\infty)$, the
function assumes opposite signs, and therefore, by the intermediate value
theorem (theorem~\ref{th:zeri}), it must vanish at least once in each of the two
intervals. On the other hand, on each of the two intervals, the function is
strictly monotonic and thus injective, so it cannot vanish more than once. We
deduce that the function $f$ vanishes at exactly two points $x_1$, $x_2$ with
\[
  0 < x_1 < 3 < x_2.
\]

The points $x_1$ and $x_2$ are therefore uniquely determined and can be
calculated, with an arbitrarily small error, using the bisection method, as we
have seen in the proof of the intermediate value theorem.
\end{proof}

\begin{example}
Consider the function
\[
  f(x) = \arctg x + \arctg\frac 1 x.
\]
We have
\[
  f'(x) = \frac{1}{1+ x^2} + \frac{1}{1 + \frac {1}{x^2}} \frac{-1}{x^2}
    = \frac {1}{1+x^2} - \frac{1}{x^2 + 1} = 0.
\]
We observe that the function $f$ is defined on $\RR\setminus \ENCLOSE{0}$, which
is not an interval but is the union of two disjoint intervals: $(-\infty, 0)
\cup (0, +\infty)$. We can then apply the monotonicity criteria separately to
the two intervals, obtaining that $f(x)$ is constant on each of the two
intervals. Thus, there will exist $c_1$ and $c_2$ such that
\[
  f(x) = \begin{cases} c_1 & \text{if $x>0$,} \\
  c_2 & \text{if $x<0$.}
  \end{cases}
\]
We can easily determine $c_1$ and $c_2$ by observing that
\begin{align*}
c_1 &= f(1) = \arctg 1 + \arctg 1 = \frac{\pi}{2} \\
c_2 &= f(-1) = \arctg (-1) + \arctg (-1) = - \frac{\pi}{2}.
\end{align*}
\end{example}
Indeed, the function $f$, despite having a zero derivative, is not constant but
only \emph{locally constant}.

\begin{example}
Consider the function $f(x) = x^3$ whose derivative is $f'(x) = 3x^2$. For every
$x\in \RR$, we have $f'(x)\ge 0$, thus we can deduce, as we already know, that
$f$ is increasing. Instead, choosing $I = [0,+\infty)$, the corresponding open
interval is $J=(0,+\infty)$. We observe that on $J$, we have $f'(x) > 0$, so we
can conclude that $f$ is strictly increasing on the entire $I$. The same holds
for the interval $(-\infty,0]$. Combining the two, we can conclude that $f(x) =
x^3$ is strictly increasing on the entire $\RR$, despite $f'(0)=0$. This shows
that a strictly monotonic function can have a zero derivative at a point.
\end{example}

More generally, it is easy to observe that if $f$ is monotonic but not strictly
monotonic, it means there are two points $a$ and $b$ for which $f(a) = f(b)$.
But if $f$ is monotonic, then for every $x\in [a,b]$, we must have $f(x) = f(a)
= f(b)$ (for example: if $f$ is increasing, we should have $f(a) \le f(x) \le
f(b)$, but if $f(a)=f(b)$, necessarily $f(x)=f(a)=f(b)$). Thus, $f$ would be
constant on $[a,b]$, and in particular, we would have an uncountable infinity of
points where the derivative is zero. This means that if $f'(x)\ge 0$ on an
interval and if $f'(x)=0$ on a finite or even countable number of points, or,
again, on a set of points with an empty interior, then $f$ is still strictly
increasing. A similar reasoning naturally applies to decreasing functions as
well.

\begin{exercise}
Prove that
\[
  \cos x \ge 1- \frac{x^2}{2} \qquad \forall x \in \RR.
\]
\end{exercise}

\begin{exercise}
Consider the function $f\colon \RR \to \RR$ defined by
\[
  f(x) = 2e^{x-1} - x^2.
\]
Show that $f$ is bijective and that the inverse function $f^{-1}\colon \RR \to
\RR$ is differentiable at all points except at the point $1$ where it has a
vertical inflection point. Calculate $(f^{-1})'(2/e)$.
\end{exercise}
\begin{proof}[Solution.]
We have
\[
  f'(x) = 2e^{x-1} - 2x, \qquad f''(x) = 2e^{x-1}-2.
\]
Thus, $f''(x) > 0$ for $x > 1$ and $f''(x)< 0$ for $x<1$, and by the
monotonicity criteria, $f'$ is strictly increasing on $[1,+\infty)$ and strictly
decreasing on $(-\infty, 1]$. Since $f'(1)=0$, it follows that $f'(x)\ge 0$ for
every $x\in \RR$ and $f'(x)=0$ only for $x=1$. Thus, $f$ is increasing by the
monotonicity criterion but also strictly increasing because if it were
increasing but not strictly, there would have to be an entire interval where
$f'$ is zero. Thus, $f$ is injective. Since $f(x)\to \pm\infty$ as $x\to
\pm\infty$, we have $\sup f(\RR) = +\infty$, $\inf f(\RR) = -\infty$, and by the
intermediate value theorem, we obtain that $f(\RR)=\RR$. Therefore, $f$ is also
surjective and hence invertible.

The inverse function of $f$ is also monotonic, thus it is continuous by
Theorem~\ref{th:monotona_continua} and is differentiable at points corresponding
to points where $f$ has a non-zero derivative. The only point where $f$ has a
zero derivative is $x=1$, and since $f(1) = 1$, it follows that $f^{-1}(y)$ is
differentiable for every $y\neq 1$ and satisfies
\[
  (f^{-1})'(f(x)) = \frac{1}{f'(x)} = \frac{1}{2 e^{x-1}-2x}.
\]
Observing that $f(0)=2/e$, we find
\[
  (f^{-1})'(2/e) = \frac{1}{f'(0)} = \frac{1}{2/e} = \frac e 2.
\]
\end{proof}